\documentclass[11pt]{article}
\usepackage{amsmath,amssymb,amsthm}
\newtheorem{proposition}{Proposition}
\begin{document}
\title{Height of a product of odd-index abscissae}
\author{}\date{}\maketitle
Put $K=\mathbb Q(\sqrt2)$, $d=3-\sqrt2$, $D=d^2$, and
$E:y^2=x^3-Dx$. Write $h$ for the absolute logarithmic Weil height.
\begin{proposition}
Let $P\in E(K)$ and let $n_1,\ldots,n_k$ be odd integers such that
$P_i=[n_i]P$ is a nonidentity point in $2E(K)$ for each $i$.
Assume, as in the square-abscissa chart of the construction, that
$x_i=x(P_i)$ is positive at both real embeddings. Then
\[
 \sum_{i=1}^k h(x_i)-(k-1)\log7
 \ \leq\ h\left(\prod_{i=1}^k x_i\right)
 \ \leq\ \sum_{i=1}^k h(x_i).
\]
In particular the lower error is $3\log7$ for four odd-index points.
\end{proposition}
\begin{proof}
Fix a finite prime $\mathfrak p$ of $K$, and normalize $v=v_{\mathfrak p}$
to have value group $\mathbb Z$. Since $D$ is integral, the points reducing
to the point at infinity on the displayed integral Weierstrass equation
form its formal subgroup
\[
 N=\{O\}\cup\{Q\in E(K_{\mathfrak p}):v(x(Q))<0\}.
\]
This is valid also at primes of bad reduction: the point at infinity is
smooth on the integral Weierstrass equation. The subgroup is open, so
$E(K_{\mathfrak p})/N$ is finite. The point $T=(0,0)$ has order two and
does not belong to $N$.

Suppose some $v(x_i)<0$. The image of $P$ in the finite quotient then has
order dividing the odd integer $n_i$, and hence has odd order. No multiple
of that image can equal the nonzero order-two image of $T$. If
$v(x_j)>v(D)$, however, the translation identity
\[
                x(P_j+T)=-D/x_j
\]
would give $P_j+T\in N$, a contradiction. Consequently, if any $v(x_i)$
is negative, every positive $v(x_j)$ is at most $v(D)$. There are at most
$k-1$ positive valuations in this case. Thus, for every finite prime,
\[
 \max\left(0,-\sum_i v(x_i)\right)
 \geq \sum_i\max(0,-v(x_i))-(k-1)v(D).
\]
If no valuation is negative the same inequality is immediate.

At a real embedding, $x_i>0$ and the curve equation gives
$x_i(x_i^2-D)\geq0$, so $x_i\geq\sqrt D=d$ in that embedding. Both
conjugates $3\pm\sqrt2$ of $d$ are greater than one. Therefore the
archimedean local height of the product is exactly the sum of the
archimedean local heights. Summing the finite inequalities with weights
$\frac12\log N\mathfrak p$, and using $N_{K/\mathbb Q}(D)=49$, gives
\[
 h\left(\prod_i x_i\right)
 \geq\sum_i h(x_i)-\frac{k-1}{2}\log49.
\]
The upper bound is the usual product inequality for Weil height.
\end{proof}

\begin{proposition}
For the construction's fixed point
\[
 P_*=(57/8-(7/2)\sqrt2,\,-165/8+(417/32)\sqrt2)
\]
and odd nonzero indices as above, one has the exact identity
\[
 h\left(\prod_i x([n_i]P_*)\right)=\sum_i h(x([n_i]P_*)).
\]
\end{proposition}
\begin{proof}
The preceding proof gives local additivity at every finite prime not
dividing $D$. The only prime dividing $D$ is $(d)$ over seven. The
Weierstrass equation is minimal there, since its discriminant has
valuation six. Reducing $\sqrt2$ to three gives
$P_*\bmod(d)=(1,6)$, a nonsingular point on $y^2=x^3$ over
$\mathbb F_7$. The nonsingular-reduction subgroup $E_0(K_{(d)})$ therefore
contains every multiple of $P_*$. If $v_{(d)}(x([n]P_*))>0$, the curve
equation would also give positive ordinate valuation, and the point
would reduce to the singular point $(0,0)$. This is impossible. Thus
all the abscissa valuations at $(d)$ are nonpositive, which again gives
local additivity. The archimedean additivity was proved above.
\end{proof}

For $U=\prod_i x_i=u_0+u_1\sqrt2$, the elementary height inequality gives
\[
 h(U)\leq h(u_0)+h(u_1)+\tfrac32\log2.
\]
Thus the same two-component comparison that bounds the height of one
elliptic abscissa can bound the sum of all four abscissa heights, with no additive loss for the fixed generator $P_*$. The product is totally positive, so
$u_0>0$ and the fixed positive scales $u_0^2$ and $(1+u_1^2)^2$ retain
their domain properties.

The proposition does not assert that arbitrary labels in the current
construction are odd. A replacement of the index encoding and its
completeness proof is required before reducing the detector family.
\end{document}
